\labsection{2}{BMI conditions and string manipulation}{sec:lab02}

\begin{labproject}{Lab 2b guided practice}
\textbf{Coverage.} Chapter~\ref{chap:logic-strings}.\\
\textbf{Goal.} Apply logical expressions and string functions to three input tasks.

\textbf{Task 1 --- BMI category flags.}

Use continuous boundaries so every valid BMI belongs to one category. The Boolean flags avoid selection statements, which are introduced in Lab 5.

\begin{lstlisting}[style=dsr]
weight <- as.numeric(readline("Enter weight (kg): "))
height <- as.numeric(readline("Enter height (m): "))

bmi <- weight / (height ^ 2)

underweight <- bmi < 18.5
normal      <- bmi >= 18.5 & bmi < 25
overweight  <- bmi >= 25 & bmi < 40
obese       <- bmi >= 40.0

print(paste("BMI:", bmi))
print(paste("Underweight:", underweight))
print(paste("Normal:", normal))
print(paste("Overweight:", overweight))
print(paste("Obese:", obese))
\end{lstlisting}

\textbf{Task 2 --- Case-insensitive comparison of two strings.}

\begin{lstlisting}[style=dsr]
string1 <- readline("Enter string 1: ")
string2 <- readline("Enter string 2: ")

same <- toupper(string1) == toupper(string2)
print(paste("Both inputs are similar:", same))
\end{lstlisting}

Converting both inputs to one case makes the comparison case-insensitive.

\textbf{Task 3 --- Uppercase a name and partially display a phone number.}

\begin{lstlisting}[style=dsr]
name <- readline("Enter name: ")
phone <- readline("Enter phone number: ")

name_upper <- toupper(name)
first3 <- substr(phone, 1, 3)
last4 <- substr(phone, nchar(phone) - 3, nchar(phone))

print(paste("Hi,", name_upper,
            ". A verification code has been sent to",
            first3, "-xxxxx", last4))
\end{lstlisting}

The output retains the first three and last four digits while masking the middle.
\end{labproject}

\begin{labdeliverable}
Submit all three programs and demonstrate a successful case-insensitive comparison.
\end{labdeliverable}
